\chapter{RUANG HILBERT}

Untuk fakta-fakta paling mendasar tentang ruang Hilbert, saya tidak dapat memikirkan pengantar yang lebih baik daripada dua bab pertama
dari buku klasik Halmos~\cite{Halmos:1951}.  Pengantar lain yang baik adalah~\cite{Young:1988}.  Bab 3 dan 9 dari~\cite{Kreyszig:1978}
serta bab 4 dari~\cite{Ha:2006} mungkin juga berguna bagi Anda.  Sebagai langkah awal menuju penjelajahan yang lebih serius
atas seluk-beluk ruang Hilbert, lihat~\cite{Conway:1990}, \cite{Conway:2000}, dan~\cite{Helmberg:1969}.


\section{Definisi dan Contoh}

Ruang Hilbert adalah \emph{ruang hasil kali dalam lengkap}.  Secara lebih terperinci:
\begin{defn} Dalam proposisi~\ref{00015025} kita menunjukkan bagaimana hasil kali dalam pada suatu ruang vektor menginduksi norma pada ruang itu dan dalam
proposisi~\ref{0001503} bagaimana norma tersebut selanjutnya menginduksi metrik.  Jika suatu ruang hasil kali dalam lengkap terhadap metrik ini,
ruang itu disebut
  \index{Hilbert!ruang}%
  \index{ruang!Hilbert}%
\df{ruang Hilbert}.  Demikian pula,
  \index{Banach!ruang}%
  \index{ruang!Banach}%
\df{ruang Banach} adalah ruang linear bernorma lengkap dan
  \index{Banach!aljabar}%
  \index{aljabar!Banach}%
\df{aljabar Banach} adalah aljabar bernorma lengkap.
\end{defn}

\begin{exam} Dengan hasil kali dalam yang didefinisikan dalam
 \index{K@$\K^n$!sebagai ruang Hilbert}%
 \index{Hilbert!ruang!$\K^n$ sebagai}%
contoh~\ref{000150012}, $\K^n$ merupakan ruang Hilbert.
\end{exam}

\begin{exam}\label{exam150013} Dengan hasil kali dalam yang diberikan dalam
 \index{l@$l_2 = l_2(\N)$!sebagai ruang Hilbert}%
 \index{l@$l_2(\Z)$!sebagai ruang Hilbert}%
 \index{Hilbert!ruang!$l_2 = l_2(\N)$, $l_2(\Z)$ sebagai}%
 \index{l@$l_2 = l_2(\N)$!barisan yang kuadratnya dapat dijumlahkan}%
 \index{l@$l_2 = l_2(\N)$!sebagai ruang hasil kali dalam}%
 \index{hasil kali dalam!ruang!$l_2$ sebagai}%
contoh~\ref{000150013}, $l_2 = l_2(\N)$ dan $l_2(\Z)$ merupakan ruang Hilbert.
\end{exam}

\begin{exam} Dengan hasil kali dalam yang didefinisikan dalam contoh~\ref{000150014}, $\fml C([a,b])$ \emph{bukan} ruang Hilbert.
\end{exam}

\begin{exam}\label{exam_CX} Misalkan $\fml C(X) = \fml C(X,\C)$ adalah keluarga semua fungsi kontinu bernilai kompleks pada ruang
Hausdorff kompak~$X$.  Dengan
 \index{seragam!norma!pada $\fml C(X)$}%
 \index{norma!seragam}%
 \index{C@$\fml C(X)$!sebagai ruang Banach}%
 \index{C@$\fml C(X)$!fungsi kontinu bernilai kompleks pada~$X$}%
\df{norma seragam}
   \[ \norm f_u := \sup\{\abs{f(x)}\colon x \in X\} \]
$\fml C(X)$ merupakan ruang Banach.  Jika ruang $X$ tidak kompak, dengan cara yang sama kita dapat menjadikan $\fml C_b(X)$, yakni himpunan
 \index{C@$\fml C_b(X)$!sebagai ruang Banach}%
semua fungsi kontinu bernilai kompleks yang \emph{terbatas} pada $X$ sebagai ruang Banach.
Jika medan skalarnya adalah $\R$, kita akan
 \index{C@$\fml C(X,\R)$!fungsi kontinu bernilai real pada~$X$}%
 \index{C@$\fml C(X,\R)$!sebagai ruang Banach}%
menuliskan $\fml C(X,\R)$ untuk keluarga fungsi kontinu bernilai real pada~$X$.  Keluarga ini juga merupakan ruang Banach.
\end{exam}

\begin{exam}\label{000213} Ruang hasil kali dalam $l_c$ yang terdiri atas semua barisan $(a_n)$ bilangan kompleks yang pada akhirnya bernilai nol
(lihat contoh~\ref{exam13122}) bukan ruang Hilbert.
\end{exam}

\begin{exam}\label{00022} Misalkan $\mu$ adalah ukuran positif pada suatu $\sigma$-aljabar $\sfml A$ yang terdiri atas
subhimpunan-subhimpunan dari himpunan~$S$.  Jika Anda belum mengenal teori ukuran umum, pengantar singkat tentang ukuran Lebesgue
pada garis real seharusnya memadai untuk contoh-contoh yang muncul dalam catatan ini.  Suatu fungsi bernilai kompleks $f$ pada $S$ disebut
 \index{terukur}%
\df{terukur} jika prapeta oleh $f$ dari setiap himpunan Borel (secara ekuivalen, setiap himpunan
terbuka) di $\C$ termasuk dalam~$\sfml A$. Kita mendefinisikan relasi ekuivalensi $\sim$ pada keluarga
fungsi terukur bernilai kompleks dengan menetapkan $f \sim g$ apabila $f$ dan $g$ berbeda pada suatu himpunan
berukuran nol, yaitu apabila $\mu(\{x \in S\colon f(x) \ne g(x)\}) = 0$.  Kita menggunakan
notasi konvensional dan menyatakan kelas ekuivalensi yang memuat $f$ dengan $f$ itu sendiri (alih-alih
dengan sesuatu yang lebih logis seperti~$[f]$).  Keluarga (kelas-kelas ekuivalensi dari) fungsi
terukur bernilai kompleks pada $S$ dinyatakan dengan $\fml M(S,\mu)$. Suatu fungsi $f \in \fml M(S,\mu)$ disebut
 \index{terintegralkan kuadrat}%
 \index{M@$\fml M(S,\mu)$ (kelas ekuivalensi fungsi terukur)}%
 \index{L@$\fml L_2(S,\mu)$!fungsi terintegralkan kuadrat}%
\df{terintegralkan kuadrat} jika $\int_S \abs{f(x)}^2\,d\mu(x) < \infty$.  Keluarga
(kelas-kelas ekuivalensi dari) fungsi terintegralkan kuadrat pada $S$ dinyatakan dengan~$\lfs 2(S,\mu)$.
 \index{L@$\lfs 2(S,\mu)$!kelas fungsi terintegralkan kuadrat}%
 \index{L@$\lfs 2(S,\mu)$!sebagai ruang Hilbert}%
 \index{Hilbert!ruang!$\lfs 2(S,\mu)$ sebagai}%
Untuk setiap $f$, $g \in \lfs 2(S,\mu)$, definisikan
   \[ \langle f,g \rangle := \int_S f(x) \conj{g(x)}\,d\mu(x)\,. \]
Dengan hasil kali dalam ini (dan operasi ruang vektor titik demi titik yang sudah jelas), $\lfs 2(S,\mu)$ merupakan ruang Hilbert.
\end{exam}

\begin{exam}\label{exam_L1S} Gunakan notasi seperti dalam contoh sebelumnya.  Suatu fungsi $f \in \fml M(S,\mu)$ disebut
 \index{terintegralkan}%
 \index{L@$\fml L_1(S,\mu)$!fungsi terintegralkan}%
\df{terintegralkan} jika $\int_S \abs{f(x)}\,d\mu(x) < \infty$.  Keluarga
(kelas-kelas ekuivalensi dari) fungsi terintegralkan pada $S$ dinyatakan dengan~$\lfs 1(S,\mu)$.
 \index{L@$\lfs 1(S,\mu)$!kelas fungsi terintegralkan}%
 \index{L@$\lfs 1(S,\mu)$!sebagai ruang Banach}%
 \index{Banach!ruang!$\lfs 1(S,\mu)$ sebagai}%
Untuk setiap $f \in \lfs 1(S,\mu)$, definisikan
    \[ {\norm f}_1 := \int_S \abs{f(x)}\,d\mu(x)\,. \]
Dengan norma ini (dan operasi ruang vektor titik demi titik yang sudah jelas), $\lfs 1(S,\mu)$ merupakan ruang Banach.
\end{exam}

\begin{defn} Jika $J$ adalah sembarang interval di $\R$, suatu fungsi bernilai real (atau kompleks) $f$ pada $J$ disebut
 \index{mutlak!kekontinuan!fungsi}%
 \index{kekontinuan!mutlak!fungsi}%
 \index{kontinu!mutlak}%
\df{kontinu mutlak} jika memenuhi syarat berikut:
  \quote untuk setiap $\epsilon > 0$ terdapat $\delta > 0$ sedemikian
 sehingga jika $(a_1,b_1), \dots,(a_n,b_n)$ \\
 merupakan subinterval-subinterval tak kosong yang saling lepas dari $J$ dengan $\sum_{k=1}^n(b_k - a_k) < \delta$, \\
 maka $\sum_{k=1}^n\abs{f(b_k) - f(a_k)}~<~\epsilon$.
  \endquote
\end{defn}

\begin{exam}\label{exam_Hsp_abs_cont} Misalkan $\lambda$ adalah ukuran Lebesgue pada interval $[0,1]$ dan $H$ adalah himpunan semua
fungsi kontinu mutlak pada $[0,1]$ sedemikian sehingga $f\,'$ termasuk dalam $\lfs 2([0,1],\lambda)$ dan $f(0) = 0$.
 \index{Hilbert!ruang!fungsi kontinu mutlak}%
 \index{mutlak!fungsi kontinu!ruang Hilbert}%
Untuk $f$ dan $g$ dalam $H$, definisikan
   \[ \langle f,g \rangle = \int_0^1 f\,'(t)\clo{g\,'(t)}\,dt. \]
Ini merupakan hasil kali dalam pada $H$ yang menjadikan $H$ ruang Hilbert.
\end{exam}

\begin{exam} Jika $H$ dan $K$ merupakan ruang Hilbert, jumlah langsung (ortogonal eksternal) keduanya $H \oplus K$
 \index{Hilbert!ruang!jumlah langsung ruang Hilbert sebagai}%
 \index{langsung!jumlah!ruang Hilbert}%
(sebagaimana didefinisikan dalam~\ref{0001504}) juga merupakan ruang Hilbert.
\end{exam}


















\section{Penjumlahan Tak Berurutan}

Mendefinisikan \emph{penjumlahan tak berurutan} dalam ruang linear bernorma tidak lebih sulit daripada dalam~$\R$.  Kita memperumum
definisi~\ref{def_summable_R}.

\begin{defn} Misalkan $V$ adalah ruang linear bernorma dan $A \subseteq V$.
Untuk setiap $F \in \fin A$, definisikan
  \[ \sbsb sF = \sum F\]
di mana $\sum F$ adalah jumlah dari vektor-vektor (yang berhingga banyaknya) dalam~$F$.
Dengan demikian, $s = \bigl(\sbsb sF\bigr)_{F \in \fin A}$ adalah jaring dalam~$V$. Jika jaring ini konvergen, himpunan
$A$ disebut
 \index{dapat dijumlahkan!himpunan vektor}%
\df{dapat dijumlahkan}; limit jaring tersebut adalah
 \index{jumlah!himpunan yang dapat dijumlahkan}%
 \index{<sum@$\sum A$ (jumlah suatu himpunan vektor)}%
\df{jumlah} dari $A$ dan dinyatakan dengan~$\sum A$. Jika $\{\norm a\colon a \in A\}$ dapat dijumlahkan, kita mengatakan bahwa himpunan $A$
 \index{mutlak!dapat dijumlahkan}%
 \index{dapat dijumlahkan!secara mutlak}%
\df{dapat dijumlahkan secara mutlak}.

Untuk mengakomodasi jumlah yang memuat vektor yang sama lebih dari satu kali, kita perlu menggunakan keluarga vektor terindeks.  Keluarga demikian
memerlukan pendekatan yang sedikit berbeda. Misalkan, sebagai contoh, $A = \bigl(x_i\bigr)_{i \in I}$ dengan $I$
sebagai sembarang himpunan indeks. Maka untuk setiap $F \in \fin I$,
  \[ \sbsb sF = \sum_{i \in F}x_i . \]
Seperti di atas, $s$ adalah jaring dalam~$V$. Jika jaring ini konvergen, $\bigl(x_i\bigr)_{i \in I}$ merupakan
 \index{dapat dijumlahkan!keluarga vektor terindeks}%
keluarga yang dapat dijumlahkan dan limitnya adalah
 \index{jumlah!keluarga terindeks}%
 \index{<sum@$\sum_{i \in I} x_i$ (jumlah suatu keluarga terindeks)}%
jumlah keluarga terindeks tersebut, yang dinyatakan dengan $\sum_{i \in I} x_i$.  Cara lain untuk mengatakan bahwa $\bigl(x_i\bigr)_{i \in I}$
dapat dijumlahkan adalah dengan mengatakan bahwa ``deret'' $\sum_{i \in I}x_i$
 \index{konvergensi!deret tak hingga}%
 \index{jumlah!deret tak hingga}%
\df{konvergen} atau bahwa ``jumlah'' $\sum_{i \in I}x_i$ \df{ada}.  Suatu keluarga terindeks $\bigl(x_i\bigr)_{i \in I}$
 \index{mutlak!dapat dijumlahkan}%
 \index{dapat dijumlahkan!secara mutlak}%
\df{dapat dijumlahkan secara mutlak} jika keluarga terindeks $\bigl(\norm{x_i}\bigr)_{i \in I}$ dapat dijumlahkan.  Cara lain untuk mengatakan
bahwa $\bigl(x_i\bigr)_{i \in I}$ dapat dijumlahkan secara mutlak adalah dengan mengatakan bahwa ``deret'' $\sum_{i \in I}x_i$
 \index{konvergensi!mutlak}%
 \index{mutlak!konvergensi}%
\df{konvergen mutlak}.

Seperti biasa, barisan diperlakukan sebagai kasus tersendiri.  Misalkan $(a_k)$ adalah barisan dalam~$V$.  Jika deret tak hingga
$\sum_{k=1}^\infty a_k$ (yaitu, barisan
 \index{parsial!jumlah}%
 \index{jumlah!parsial}%
 \index{deret tak hingga!jumlah parsial}%
 \index{deret!tak hingga}%
jumlah parsial $s_n = \sum_{k=1}^n a_k$) konvergen ke suatu
vektor $b$ dalam $V$, kita mengatakan bahwa barisan $(a_k)$
 \index{deret tak hingga!dalam ruang linear bernorma}%
 \index{dapat dijumlahkan!barisan}%
 \index{barisan!dapat dijumlahkan}%
\df{dapat dijumlahkan} atau, secara ekuivalen, bahwa deret $\sum_{k=1}^\infty a_k$ merupakan
 \index{konvergensi!deret}%
 \index{deret!konvergen}%
\df{deret konvergen}.  Vektor $b$ disebut
 \index{jumlah!deret tak hingga}%
 \index{deret!jumlah}%
 \index{deret tak hingga!jumlah}%
\df{jumlah} deret $\sum_{k=1}^\infty a_k$ dan kita menulis
  \[ \sum_{k=1}^\infty a_k = b\,. \]
Jelas bahwa syarat perlu dan cukup agar suatu deret $\sum_{k=1}^\infty a_k$
konvergen atau, secara ekuivalen, agar barisan $(a_k)$ dapat dijumlahkan, adalah adanya suatu
vektor $b$ dalam $V$ sedemikian sehingga
  \begin{equation}\label{cond_ser_conv}
     \biggnorm{b - \sum_{k=1}^n a_k} \sto 0 \qquad \text{ketika} \qquad n \sto \infty.
  \end{equation}
\end{defn}

\begin{prop}\label{prop_abs_sum_implies_sum} Suatu ruang linear bernorma $V$ lengkap jika dan hanya jika setiap barisan yang
dapat dijumlahkan secara mutlak dalam $V$ dapat dijumlahkan.
\end{prop}

\begin{proof}[\emph{Petunjuk pembuktian}] Misalkan $(x_n)$ adalah barisan Cauchy dalam suatu ruang linear bernorma tempat setiap barisan yang
dapat dijumlahkan secara mutlak juga dapat dijumlahkan.  Carilah suatu subbarisan $\bigl(x_{n_k}\bigr)$ dari $(x_n)$ sedemikian sehingga $\norm{x_n - x_m} < 2^{-k}$
apabila $m$, $n \ge n_k$.  Jika kita menetapkan $y_1 = x_{n_1}$ dan $y_j = x_{n_j} - x_{n_{j-1}}$ untuk $j > 1$, maka $(y_j)$
dapat dijumlahkan secara mutlak.   \ns
\end{proof}

\begin{prop} Jika $\bigl(x_i\bigr)_{i \in I}$ dan $\bigl(y_i\bigr)_{i \in I}$ merupakan
keluarga terindeks yang dapat dijumlahkan dalam suatu ruang linear bernorma dan $\alpha$ adalah skalar, maka
$\bigl(\alpha x_i\bigr)_{i \in I}$ dan $\bigl(x_i + y_i\bigr)_{i \in I}$
dapat dijumlahkan; selain itu,
  \[ \sum_{i \in I} \alpha x_i = \alpha \, \sum_{i \in I} x_i \]
dan
  \[ \sum_{i \in I} x_i + \sum_{i \in I} y_i =
                         \sum_{i \in I} (x_i + y_i). \]
\end{prop}

\begin{rem} Dalam buku teks sering dijumpai suatu versi dari
``proposisi'' berikut.
\vskip 10 pt
\begin{quote}
Jika $\sum_{i \in I}x_i = u$ dan $\sum_{i \in I}x_i = v$ dalam suatu
ruang Hilbert (atau Banach), maka $u = v$.
\end{quote}

 \vskip 10 pt
\noindent (Saya pernah melihat bukti sepanjang 6 baris untuk pernyataan ini.) Pernyataan ini merupakan
konsekuensi sepele dari hasil yang mana?
\end{rem}

\begin{prop} Jika $\bigl(x_i\bigr)_{i \in I}$ adalah keluarga terindeks vektor yang dapat dijumlahkan dalam ruang Hilbert $H$ dan
$y$ adalah vektor dalam $H$, maka $\bigl(\langle x_i,y \rangle\bigr)_{i \in I}$ adalah keluarga terindeks skalar yang dapat dijumlahkan dan
   \[ \sum_{i \in I}\langle x_i,y \rangle  = \biggl < \sum_{i \in I} x_i, y \biggr >. \]
\end{prop}

Berikut ini merupakan perumuman dari contoh~\ref{exam150013}.

\begin{exam} Misalkan $I$ adalah sembarang himpunan tak kosong dan $l^2(I)$ adalah himpunan semua fungsi bernilai kompleks $x$ pada $I$ sedemikian sehingga
$\sum_{i \in I}\abs{x_i}^2 < \infty$. Untuk semua $x$, $y \in l^2(I)$, definisikan
    \[ \langle x,y \rangle = \sum_{i \in I}x(i)\conj{y(i)}.\]
  \begin{enumerate}
    \item Jika $x \in l^2(I)$, maka $x(i) = 0$ untuk semua indeks $i$, kecuali mungkin indeks-indeks dalam suatu himpunan terhitung.
    \item Pemetaan $(x,y) \mapsto \langle x,y \rangle$ merupakan hasil kali dalam pada ruang vektor $l^2(I)$.
    \item Ruang $l^2(I)$ merupakan ruang Hilbert.
  \end{enumerate}
\end{exam}

\begin{prop}\label{prop_sum_conv} Misalkan $(x_n)$ adalah barisan dalam suatu ruang Hilbert.  Jika barisan tersebut, ketika dipandang sebagai keluarga terindeks,
dapat dijumlahkan, maka deret tak hingga $\sum_{k=1}^\infty x_k$ konvergen.
\end{prop}

\begin{exam} Kebalikan dari proposisi~\ref{prop_sum_conv} tidak berlaku.
\end{exam}

Jumlah langsung dari dua ruang Hilbert merupakan ruang Hilbert.

\begin{prop} Jika $H$ dan $K$ merupakan ruang Hilbert, maka jumlah langsung (ortogonal) keduanya (sebagaimana didefinisikan dalam~\ref{0001504})
merupakan ruang Hilbert.
\end{prop}

Kita juga dapat mengambil hasil kali lebih dari dua ruang Hilbert---bahkan hasil kali suatu keluarga tak hingga ruang Hilbert.

\begin{prop} Misalkan $\bigl(H_i\bigr)_{i \in I}$ adalah keluarga terindeks ruang Hilbert.  Suatu fungsi $x\colon I \sto \bigcup H_i$
termasuk dalam $H = \D\bigoplus_{i \in I}H_i$ jika $x(i) \in H_i$ untuk setiap $i$ dan jika $\sum_{i \in I}\norm{x(i)}^2 < \infty$.
Mendefinisikan $\langle \hphantom x, \hphantom y \rangle$ pada $H$ dengan $\langle x, y \rangle := \sum_{i \in I} \langle
x(i), y(i) \rangle$ menjadikan $H$ ruang Hilbert.
\end{prop}

\begin{defn} Ruang Hilbert yang didefinisikan dalam proposisi sebelumnya adalah
 \index{eksternal!jumlah langsung ortogonal}%
 \index{ortogonal!jumlah langsung!ruang Hilbert}%
 \index{langsung!jumlah!ruang Hilbert}%
 \index{Hilbert!ruang!jumlah langsung}%
 \index{jumlah!langsung}%
\df{jumlah langsung ortogonal (eksternal)} dari ruang-ruang Hilbert~$H_i$.
\end{defn}

\begin{exer}  Apakah jumlah langsung yang didefinisikan di atas (dengan morfisme yang sesuai) merupakan hasil kali dalam kategori $\cat{HSp}$
yang objeknya ruang Hilbert dan morfismenya pemetaan linear terbatas?  Apakah jumlah langsung itu merupakan koproduk?
\end{exer}


























\section{Geometri Ruang Hilbert}

\begin{conv}  Dalam konteks ruang Hilbert, kata ``subruang'' akan selalu berarti
 \index{konvensi!subruang dari ruang Hilbert bersifat tertutup}%
 \index{subruang!ruang Hilbert}%
\emph{subruang vektor tertutup}.  Alasannya adalah bahwa kita ingin agar subruang dari suatu ruang Hilbert merupakan ruang Hilbert tersendiri; dengan kata lain,
subruang dari suatu ruang Hilbert seharusnya merupakan subobjek dalam kategori ruang Hilbert (dan morfisme yang sesuai).  Untuk menunjukkan bahwa $M$ adalah subruang dari $H$, kita menulis
 \index{<binrelle@$\preccurlyeq$ (subruang dari ruang Banach atau Hilbert)}%
$M \preccurlyeq H$.  Subruang vektor (yang tidak harus tertutup) dari suatu ruang Hilbert juga
 \index{linear!subruang}%
 \index{vektor!subruang}%
 \index{linear!manifold}%
 \index{manifold!linear}%
disebut \emph{subruang linear} atau \emph{manifold linear}.
\end{conv}

\begin{defn} Misalkan $A$ adalah subhimpunan tak kosong dari ruang Banach (atau Hilbert)~$B$.  Kita mendefinisikan
 \index{tertutup!rentang linear}%
 \index{linear!rentang!tertutup}%
 \index{<unops@$\bigvee A$ (rentang linear tertutup dari~$A$)}%
\df{rentang linear tertutup} dari $A$ (dinyatakan dengan $\bigvee A$) sebagai subruang terkecil dari $B$ yang memuat~$A$; yaitu,
irisan semua subruang dari $B$ yang memuat~$A$.
\end{defn}

Seperti pada definisi \ref{smplx005def}, \ref{def_interior}, dan \ref{def_closure}, kita perlu menunjukkan bahwa definisi sebelumnya bermakna.

\begin{prop} Definisi sebelumnya bermakna.  Lebih lanjut, suatu vektor berada dalam rentang linear tertutup dari $A$ jika dan hanya jika vektor itu termasuk
dalam penutupan rentang linear dari~$A$.
\end{prop}

\begin{thm}[Teorema Vektor Peminimum] Jika $C$ adalah subhimpunan konveks tertutup tak kosong dari suatu
 \index{teorema vektor peminimum}%
ruang Hilbert $H$ dan $a \in C^c$, maka terdapat tepat satu $b \in C$ sedemikian sehingga
$\norm{b - a} \le \norm{x - a}$ untuk setiap $x \in C$.
 \[\xymatrix{
     &&&&&{}  \\
     &&&& x   \\
     a\ar[rrrru]\ar[rrr] &&& b  \\
     &&&&&C   \\
     &&&&{}\ar@{-}@/^3.5pc/[uuuu]       }\]
\end{thm}

\begin{proof}[\emph{Petunjuk pembuktian}]  Untuk bagian eksistensi, misalkan $d = \inf\{\norm{x - a}\colon x \in C\}$.  Pilih suatu barisan
$(y_n)$ vektor dalam $C$ sedemikian sehingga $\norm{y_n - a} \sto d$.  Untuk membuktikan bahwa barisan $(y_n)$ adalah barisan Cauchy, terapkan
\emph{hukum jajaran genjang}~\ref{parallelogram_law} pada vektor $y_m - a$ dan $y_n - a$ untuk $m$, $n \in \N$.  Bagian keunikan
serupa: terapkan \emph{hukum jajaran genjang} pada vektor $b - a$ dan $c - a$, dengan $b$ dan $c$ sebagai dua vektor yang
memenuhi kesimpulan teorema.  \ns
\end{proof}

\begin{exam} Ruang vektor $\R^2$ dengan metrik seragam merupakan ruang Banach. Untuk melihat
bahwa \emph{Teorema Vektor Peminimum} tidak berlaku dalam ruang ini, ambillah $C$ sebagai
bola satuan tertutup berpusat di titik asal dan $a$ sebagai titik~$(2,0)$.
\end{exam}

\begin{exam} Himpunan-himpunan
   \[ C_1 = \biggl\{f \in \fml C([0,1])\colon \int_0^{1/2}f - \int_{1/2}^1f = 1\biggr\} \qquad \text{dan} \qquad
                               C_2 = \biggl\{f \in \lfs 1([0,1],\lambda)\colon \int_0^1f = 1\biggr\} \]
merupakan contoh yang menunjukkan bahwa baik pernyataan eksistensi maupun pernyataan keunikan dalam \emph{Teorema Vektor Peminimum}
tidak harus berlaku dalam ruang Banach. (Simbol $\lambda$ menyatakan ukuran Lebesgue.)
\end{exam}

\begin{thm}[Teorema Dekomposisi Vektor]\label{thm_vec_decomp} Misalkan $H$ adalah ruang Hilbert dan
 \index{vektor!teorema dekomposisi}%
$M$ adalah subruang (linear tertutup) dari~$H$.
Maka untuk setiap $x \in H$ terdapat tepat satu pasangan vektor $y \in M$ dan $z \in M^\perp$ sedemikian sehingga $x = y + z$.
\end{thm}

\begin{proof}[\emph{Petunjuk pembuktian}] Asumsikan $x \notin M$ (jika tidak, hasilnya sepele). Jelaskan bagaimana kita
mengetahui bahwa terdapat tepat satu vektor $y \in M$ sedemikian sehingga
   \begin{equation} \norm{x - y} \le \norm{x - m} \end{equation}
untuk semua $m \in M$.  Jelaskan mengapa, untuk bagian eksistensi dari bukti, cukup dibuktikan
bahwa $x - y$ tegak lurus terhadap setiap vektor \emph{satuan} $m \in M$.  Jelaskan mengapa benar
bahwa untuk setiap vektor satuan $m \in M$,
   \begin{equation} \norm{x - y}^2 \le \norm{x - (y + \lambda m)}^2 \end{equation}
di mana $\lambda := \langle x - y, m \rangle$.  \ns
\end{proof}

\begin{cor} Jika $M$ merupakan subruang dari suatu ruang Hilbert~$H$, maka $H = M \oplus M^\perp$.
\end{cor}

\begin{exam} Hasil sebelumnya menyatakan bahwa setiap subruang dari suatu ruang Hilbert merupakan suku
langsung. Hasil ini tidak harus benar jika $M$ hanya diasumsikan sebagai subruang linear dari ruang tersebut.
Sebagai contoh, perhatikan bahwa $M = l_c$ (lihat contoh~\ref{000213}) adalah subruang linear dari
ruang Hilbert $l_2$ (lihat contoh~\ref{exam150013}), tetapi $M^\perp = \{\vc 0\}$.
\end{exam}

\begin{cor} Suatu subruang vektor $A$ dari ruang Hilbert $H$ padat dalam $H$ jika dan hanya jika $A^\perp = \{0\}$.
\end{cor}

\begin{prop} Setiap subruang sejati dari suatu ruang Hilbert memiliki interior kosong.
\end{prop}

\begin{prop} Jika suatu basis Hamel bagi ruang Hilbert tidak hingga, maka basis tersebut tak terhitung.
\end{prop}

\begin{proof}[\emph{Petunjuk pembuktian}] Andaikan basis tak hingga terhitungnya ialah $\{e^1, e^2, e^3, \dots \}$.
Untuk setiap $n \in \N$, misalkan $M_n = \spn\{e^1, e^2, \dots, e^n\}$.  Terapkan \emph{teorema kategori Baire} (lihat~\cite{Erdman:2007},
teorema 29.3.20) pada $\bigcup_{n=1}^\infty M_n$ dengan mengingat proposisi sebelumnya.   \ns
\end{proof}

\begin{prop}\label{00023} Misalkan $H$ adalah ruang Hilbert. Maka pernyataan-pernyataan berikut berlaku:
 \begin{enumerate}
  \item jika $A \subseteq H$, maka $A \subseteq A^{\perp\perp}$;
  \item jika $A \subseteq B \subseteq H$, maka $B^\perp \subseteq A^\perp$;
  \item jika $M$ adalah subruang dari $H$, maka $M = M^{\perp\perp}$; dan
  \item jika $A \subseteq H$, maka $\bigvee A = A^{\perp\perp}$.
 \end{enumerate}
\end{prop}

\begin{prop} Misalkan $M$ dan $N$ adalah subruang dari suatu ruang Hilbert.  Maka
 \begin{enumerate}
  \item  $(M + N)^\perp = (M \cup N)^\perp =  M^\perp \cap N^\perp$, dan
  \item  $(M \cap N)^\perp = \clo{M^\perp + N^\perp}$.
 \end{enumerate}
\end{prop}

\begin{prop} Misalkan $V$ adalah ruang linear bernorma.  Terdapat hasil kali dalam pada $V$ yang menginduksi norma pada $V$
jika dan hanya jika norma tersebut memenuhi \emph{hukum jajaran genjang}~\ref{parallelogram_law}.
\end{prop}

\begin{proof}[\emph{Petunjuk pembuktian}]  Kita mengatakan bahwa suatu hasil kali dalam $\langle \,\cdot\,\,,\,\cdot\, \rangle$
pada $V$ menginduksi norma $\norm{\,\cdot\,}$ jika $\norm x = \sqrt{\langle x,x \rangle}$ berlaku untuk setiap $x$ dalam $V$.
Untuk membuktikan bahwa jika suatu norma memenuhi \emph{hukum jajaran genjang}, maka norma tersebut diinduksi oleh hasil kali dalam,
gunakan persamaan dalam \emph{identitas polarisasi}~\ref{0001507} sebagai definisi.
Buktikan terlebih dahulu bahwa $\langle y,x \rangle = \conj{\langle x,y \rangle}$ untuk semua $x$, $y\in V$ dan bahwa
$\langle x,x \rangle > 0$ apabila $x \ne 0$.  Selanjutnya, untuk sembarang $z \in V$, definisikan fungsi
$f\colon V \sto \C\colon x \mapsto \langle x,z \rangle$.  Buktikan bahwa
   \[ f(x + y) + f(x - y) = 2f(x) \tag{$\ast$} \]
untuk semua $x$, $y \in V$.  Gunakan ini untuk membuktikan bahwa $f(\alpha x) = \alpha f(x)$ untuk semua
$x \in V$ dan $\alpha \in \R$.  (Mulailah dengan $\alpha$ berupa bilangan asli.)  Kemudian
tunjukkan bahwa $f$ aditif. (Jika $u$ dan $v$ adalah sembarang anggota $V$, misalkan $x = u + v$
dan $y = u - v$.  Gunakan~$(\ast)$.)  Terakhir, buktikan bahwa $f(\alpha x) = \alpha f(x)$ untuk
$\alpha$ kompleks dengan menunjukkan bahwa $f(ix) = i\,f(x)$ untuk semua $x \in V$.  \ns
\end{proof}
















\section{Himpunan Ortonormal dan Basis}\label{sec_onbases}
\begin{defn} Suatu subhimpunan tak kosong $E$ dari ruang Hilbert bersifat
 \index{ortonormal}%
\df{ortonormal} jika $e \perp f$ untuk setiap pasangan vektor berbeda $e$ dan $f$ di $E$, dan
$\norm e = 1$ untuk setiap $e \in E$.
\end{defn}

\begin{prop} Dalam ruang hasil kali dalam, setiap himpunan ortonormal bebas linear.
\end{prop}

\begin{prop}[Pertidaksamaan Bessel]\label{C063141} Misalkan $H$ ruang Hilbert dan $E$ suatu
 \index{pertidaksamaan Bessel}%
subhimpunan ortonormal dari~$H$. Untuk setiap $x \in H$,
  \[ \sum_{e \in E} \norm{\langle x,e \rangle}^2 \le \norm x^2\,. \]
\end{prop}

\begin{cor} Jika $E$ subhimpunan ortonormal dari ruang Hilbert $H$ dan $x \in H$, maka
$\{e \in E\colon \langle x,e \rangle~\ne~0\}$ terhitung.
\end{cor}

\begin{prop}\label{C063144} Misalkan $E$ himpunan ortonormal dalam ruang Hilbert~$H$. Pernyataan-pernyataan
berikut ekuivalen.
 \begin{enumerate}
  \item $E$ adalah himpunan ortonormal maksimal.
  \item Jika $x \perp E$, maka $x = 0$ ($E$ bersifat \emph{total}).
 \index{total!himpunan ortonormal}%
  \item $\bigvee E = H$ ($E$ adalah \emph{himpunan ortonormal lengkap}).
 \index{lengkap!himpunan ortonormal}%
  \item $\D x = \sum_{e \in E} \langle x,e \rangle e$ untuk semua $x \in H$. \emph{(Ekspansi
Fourier)}
 \index{Fourier!ekspansi}%
  \item $\D \langle x,y \rangle = \sum_{e \in E} \langle x,e \rangle \langle e,y \rangle$
untuk semua $x,y \in H$. \emph{(Identitas Parseval.)}
 \index{identitas Parseval}%
  \item $\D \norm{x}^2 = \sum_{e \in E} \abs{\langle x,e \rangle}^2$ untuk semua $x \in H$.
\emph{(Identitas Parseval.)}
 \end{enumerate}
\end{prop}

\begin{defn}\label{C063147} Jika $H$ ruang Hilbert, maka himpunan ortonormal $E$ yang memenuhi
salah satu (dan karenanya semua) kondisi dalam proposisi~\ref{C063144} disebut
 \index{ortonormal!basis}%
 \index{basis!ortonormal}%
\df{basis ortonormal} untuk~$H$ (atau
 \index{lengkap!himpunan ortonormal}%
 \index{ortonormal!himpunan!lengkap}%
\df{himpunan ortonormal lengkap} dalam~$H$, atau
 \index{Hilbert!ruang!basis untuk}%
 \index{basis!untuk ruang Hilbert}%
\df{basis ruang Hilbert} untuk~$H$). Perhatikan bahwa konsep ini sangat berbeda dari
basis biasa (Hamel) untuk ruang vektor (lihat~\ref{defn_Hamel_basis}). Untuk ruang Hilbert,
vektor-vektor basis harus memiliki panjang satu dan saling tegak lurus
(kondisi yang tidak bermakna dalam ruang vektor); tetapi \emph{tidak} disyaratkan bahwa
setiap vektor dalam ruang tersebut merupakan kombinasi linear dari vektor-vektor basis.
\end{defn}

\begin{prop} Jika $E$ subhimpunan ortonormal dari ruang Hilbert $H$, maka terdapat
basis ortonormal untuk $H$ yang memuat~$E$.
\end{prop}

\begin{cor} Setiap ruang Hilbert tak nol memiliki basis.
\end{cor}

\begin{prop} Jika $E$ subhimpunan ortonormal dari ruang Hilbert $H$ dan $x \in H$, maka
$\{\langle x,e \rangle e: e \in E\}$ dapat dijumlahkan.
\end{prop}

\begin{exam}\label{C063149} Untuk setiap $n \in \N$, misalkan $\vc e^n$ adalah barisan dalam $l_2$ yang
koordinat ke-$n$-nya bernilai $1$ dan semua koordinat lainnya bernilai~$0$. Maka $\{\vc e^n\colon n \in \N\}$
adalah basis ortonormal untuk~$l_2$. Basis ini disebut
 \index{biasa!basis ortonormal untuk~$l_2$}%
 \index{ortonormal!basis!biasa (atau standar) untuk~$l_2$}%
 \index{basis!standar (atau biasa)}%
\df{biasa} (atau
 \index{standar!basis ortonormal}%
\df{standar}); basis tersebut adalah \df{basis ortonormal} untuk~$l_2$.
\end{exam}

\begin{exer} Suatu fungsi $f\colon \R \sto \C$ berbentuk
   \[ f(t) = a_0 + \sum_{k=1}^n (a_k \cos kt + b_k \sin kt) \]
dengan $a_0, \dots, a_n, b_1, \dots, b_n$ bilangan kompleks disebut
 \index{polinom trigonometri}%
 \index{polinom!trigonometri}%
\df{polinom trigonometri}.
 \begin{enumerate}
  \item Jelaskan mengapa fungsi bernilai kompleks yang $2\pi$-periodik pada garis real $\R$
sering diidentifikasi dengan fungsi bernilai kompleks pada lingkaran satuan~$\T$.
  \item Beberapa penulis menyebut \emph{polinom trigonometri} sebagai fungsi berbentuk
    \[ f(t) = \sum_{k=-n}^n c_k e^{ikt} \]
 dengan $c_{-n}, \dots, c_{-1}, c_0, c_1, \dots, c_n$ bilangan kompleks. Jelaskan dengan cermat mengapa bentuk ini persis
sama dengan definisi sebelumnya.
  \item Jelaskan alasan penggunaan istilah \emph{polinom trigonometri}.
  \item Buktikan bahwa setiap fungsi kontinu yang $2\pi$-periodik pada $\R$ merupakan limit seragam
dari suatu barisan polinom trigonometri.
 \end{enumerate}
\end{exer}

\begin{prop}\label{prop_unique_onbasis} Misalkan $E$ basis untuk ruang Hilbert $H$ dan $x \in H$.
Jika $x = \sum_{e \in E} \alpha_e e$, maka $\alpha_e = \langle x,e \rangle$ untuk setiap $e \in E$.
\end{prop}

\begin{prop} Tinjau ruang Hilbert $\lfs 2([0,2\pi],\C)$ yang terdiri atas (kelas-kelas ekuivalensi dari) fungsi
 \index{L@$\lfs 2([0,2\pi],\C)$}%
bernilai kompleks yang terintegralkan kuadrat pada $[0,2\pi]$, dengan hasil kali dalam
  \[ \langle f,g \rangle = \frac1{2\pi}\int_0^{2\pi} f(x) \conj{g(x)}\,dx. \]
Untuk setiap bilangan bulat $n$ (positif, negatif, atau nol), misalkan $\vc e^n(x) = e^{inx}$ untuk semua $x
\in [0,2\pi]$. Maka himpunan fungsi $\vc e^n$ untuk $n \in \Z$ merupakan
basis ortonormal untuk $\lfs 2([0,2\pi],\C)$.
\end{prop}

\begin{exam} Jumlah deret tak hingga $\D\sum_{k=1}^\infty \frac1{k^2}$ adalah $\dfrac{\pi^2}6$\,.
\end{exam}

\begin{proof}[\emph{Petunjuk pembuktian}] Misalkan $f(x) = x$ untuk $0 \le x \le 2\pi$. Pandang $f$ sebagai anggota
$\lfs 2([0,2\pi])$. Tentukan $\norm{f}$ dengan dua cara. \ns
\end{proof}

Proposisi berikut memungkinkan kita mendefinisikan gagasan \emph{dimensi} ruang Hilbert. (Pada ruang berdimensi
hingga, gagasan ini sama dengan konsep dalam ruang vektor.)

\begin{prop} Sebarang dua basis ortonormal dari suatu ruang Hilbert memiliki kardinalitas yang sama (artinya, keduanya
berkorespondensi satu-satu).
\end{prop}

\begin{proof} Lihat \cite{Halmos:1951}, halaman 29, Teorema 1. \ns
\end{proof}

\begin{defn}
 \index{dimensi}%
Yang dimaksud dengan \df{dimensi} ruang Hilbert adalah kardinalitas dari sebarang basis ortonormal untuk ruang tersebut.
\end{defn}

\begin{prop}\label{prop_Schauder_basis_Hsp} Ruang Hilbert separabel jika dan hanya jika ruang tersebut memiliki
 \index{separabel!ruang Hilbert}%
 \index{Hilbert!ruang!separabel}%
basis ortonormal terhitung.
\end{prop}

\begin{proof} Lihat \cite{Kreyszig:1978}, halaman 171, Teorema 3.6-4. \ns
\end{proof}





















\section{Teorema Riesz--Fr\'echet}
\begin{exam}\label{000341} Misalkan $H$ ruang Hilbert dan $a \in H$. Definisikan $\psi_a\colon
H \sto \C\colon x \mapsto \langle x,a \rangle$. Maka $\psi_a \in H^*$.
\end{exam}

Sekarang kita menggeneralisasi teorema~\ref{0001601} ke ranah berdimensi tak hingga. Hasil ini
kadang-kadang disebut \emph{Teorema Representasi Riesz} (yang dapat menimbulkan kebingungan dengan
 \index{representasi!fungsional linear kontinu}%
 \index{representasi!teorema kecil Riesz}%
hasil yang lebih substansial mengenai representasi fungsional linear tertentu sebagai ukuran)
atau \emph{Teorema Representasi Riesz Kecil}. Teorema ini menyatakan bahwa \emph{satu-satunya}
fungsional linear kontinu pada ruang Hilbert adalah fungsi $\psi_a$ yang didefinisikan dalam~\ref{000341}.
Dengan kata lain, jika diberikan fungsional linear kontinu $f$ pada ruang Hilbert, terdapat vektor unik
$a$ dalam ruang tersebut sehingga nilai $f$ pada vektor $x$ dapat dinyatakan hanya dengan mengambil
hasil kali dalam $x$ dengan~$a$.

\begin{thm}[Teorema Riesz--Fr\'echet]\label{000342} Jika $f \in H^*$ dengan $H$ ruang Hilbert,
 \index{teorema Riesz--Fr\'echet}%
maka terdapat vektor unik $a$ dalam $H$ sehingga $f = \psi_a$. Selain itu,
$\norm a = \norm{\psi_a}$.
\end{thm}

\begin{proof}[\emph{Petunjuk pembuktian}] Untuk kasus $f \ne 0$, pilih vektor satuan $z$ dalam $(\ker f)^\perp$.
Perhatikan bahwa untuk setiap $x \in H$, vektor $x - \frac{f(x)}{f(z)}\,z$ termasuk dalam $\ker f$. \ns
\end{proof}

\begin{cor} Kernel setiap fungsional linear terbatas tak nol pada ruang Hilbert memiliki kodimensi~$1$.
\end{cor}

\begin{proof}[\emph{Petunjuk pembuktian}]
 \index{kodimensi}%
Yang dimaksud dengan \df{kodimensi} suatu subruang dari ruang Hilbert adalah dimensi komplemen ortogonalnya. Pernyataan ini sebenarnya
merupakan akibat dari \emph{pembuktian} \emph{Teorema Riesz--Fr\'echet}. Dalam pembuktian tersebut kita mendapati bahwa vektor
yang merepresentasikan fungsional linear terbatas $f$ merupakan kelipatan (tak nol) dari vektor satuan dalam $(\ker f)^\perp$.
Apa yang akan terjadi jika terdapat dua vektor bebas linear dalam $(\ker f)^\perp$? \ns
\end{proof}

\begin{defn} Ingat bahwa pemetaan $T\colon V \sto W$ antara ruang-ruang vektor kompleks bersifat
 \index{linear konjugat}%
 \index{linear!konjugat}%
\df{linear konjugat} jika $T(x + y) = Tx + Ty$ dan $T(\alpha x) = \conj{\alpha}\,Tx$ untuk semua
$x$, $y \in V$ dan $\alpha \in \C$. Pemetaan linear konjugat yang bijektif disebut
 \index{antiisomorfisme}%
 \index{isomorfisme!anti-}%
\df{antiisomorfisme}.
\end{defn}

\begin{exam}\label{0022057} Untuk ruang Hilbert $H$, pemetaan
$\psi\colon H \sto H^*\colon a \mapsto \psi_a$ yang didefinisikan dalam contoh~\ref{000341} merupakan
antiisomorfisme isometrik antara $H$ dan ruang dualnya.
\end{exam}

\begin{defn} Pemetaan linear konjugat $C\colon V \sto V$ dari ruang vektor ke dirinya sendiri
yang memenuhi $C^2 = \id V$ disebut
 \index{konjugasi}
\df{konjugasi} pada~$V$.
\end{defn}

\begin{exam} Konjugasi kompleks $z \mapsto \conj z$ merupakan contoh konjugasi dalam
ruang vektor~$\C$.
\end{exam}

\begin{exam}\label{0022067} Misalkan $H$ ruang Hilbert dan $E$ basis untuk~$H$. Maka
pemetaan $x \mapsto \sum_{e \in E}\langle e,x \rangle e$ merupakan konjugasi dan isometri pada~$H$.
\end{exam}

Istilah ``antiisomorfisme'' agak menyesatkan. Istilah tersebut mengesankan sesuatu yang sama sekali
berbeda dari isomorfisme. Padahal, antiisomorfisme hampir sebaik isomorfisme.
Pada intinya, proposisi berikut menyatakan bahwa ruang Hilbert dan ruang dualnya isomorfik
\emph{karena} keduanya antiisomorfik.

\begin{prop}\label{0022071} Misalkan $H$ ruang Hilbert, $\psi\colon H \sto H^*$ adalah
antiisomorfisme yang didefinisikan dalam~\ref{000341} (lihat~\ref{0022057}), dan $C$ adalah konjugasi
yang didefinisikan dalam~\ref{0022067}. Maka komposisi $C \psi^{-1}$ merupakan isomorfisme isometrik
dari $H^*$ ke~$H$.
\end{prop}

\begin{cor}\label{cor_Hsp_iso_dual} Setiap ruang Hilbert isomorfik secara isometrik dengan ruang dualnya.
\end{cor}

\begin{prop} Misalkan $H$ ruang Hilbert dan $\psi\colon H \sto H^*$ antiisomorfisme
yang didefinisikan dalam~\ref{000341} (lihat~\ref{0022057}). Jika kita mendefinisikan $\langle f,g \rangle :=
\langle \psi^{-1}g, \psi^{-1}f \rangle$ untuk semua $f$, $g \in H^*$, maka $H^*$ menjadi
ruang Hilbert yang isomorfik secara isometrik dengan~$H$. Norma yang dihasilkan pada $H^*$ adalah norma
yang biasa digunakan.
\end{prop}

\begin{cor} Setiap ruang Hilbert bersifat refleksif.
\end{cor}

\begin{exam} Misalkan $H$ himpunan semua fungsi bernilai kompleks $f$ yang kontinu mutlak
pada $[0,1]$ sedemikian sehingga $f(0) = 0$ dan $f' \in
\lfs 2([0,1],\lambda)$. Definisikan hasil kali dalam pada $H$ dengan
 \[\langle f,g \rangle := \int_0^1 f'(t)\,\clo{g'(t)}\,dt.\]
Kita telah mengetahui bahwa $H$ ruang Hilbert (lihat contoh~\ref{exam_Hsp_abs_cont}).
Tetapkan $t \in (0,1]$. Definisikan
   \[ E_t\colon H \sto \C\colon f \mapsto f(t). \]
Maka $E_t$ merupakan fungsional linear terbatas pada~$H$. Berapakah $\norm{E_t}$?
Vektor $g$ manakah dalam $H$ yang merepresentasikan fungsional~$E_t$?
\end{exam}

\begin{exer} Pada suatu sore, kawan Anda, Fred R. Dimm, datang membawa masalah. ``Begini,'' katanya,
``Pada $\lfs 2 = \lfs 2(\R,\R)$, fungsional evaluasi $E_0$, yang mengevaluasi setiap anggota $\lfs 2$
di~$0$, jelas merupakan fungsional linear terbatas. Jadi, menurut Teorema Representasi Riesz,
seharusnya ada fungsi $g$ dalam $\lfs 2$ sedemikian sehingga $f(0) = \langle f,g \rangle = \int_{-\infty}^\infty fg$
untuk semua $f$ dalam~$\lfs 2$. Namun, itu adalah sifat `fungsi~$\delta$', yang katanya
tidak ada. Di mana letak kesalahan saya?'' Berilah Freddy nasihat (yang membantu).
\end{exer}

\begin{exer} Masih pada sore yang sama, Freddy kembali. Setelah mempertimbangkan
nasihat yang Anda berikan (dalam soal sebelumnya), ia merevisi
pertanyaannya dengan menjadikan domain $E_0$ sebagai himpunan fungsi bernilai real yang terbatas dan
kontinu pada~$\R$. Dengan sabar Anda
menjelaskan kepada Freddy bahwa sekarang setidaknya ada dua kekeliruan dalam caranya
menggunakan `fungsi~$\delta$'. Apakah kedua kekeliruan tersebut?
\end{exer}

\begin{exer} Mimpi buruk itu berlanjut. Sekarang pukul 23.00 dan Freddy kembali lagi.
Kali ini (setelah, untungnya, menyerah soal fungsi~$\delta$) ia ingin
menerapkan teorema representasi pada fungsional ``evaluasi
pada nol'' pada ruang~$l_2(\Z)$. Dapatkah ia melakukannya? Jelaskan.
\end{exer}



























\section{Topologi Kuat dan Lemah pada Ruang Hilbert}
Kita mengingat kembali secara singkat definisi dan fakta-fakta penting mengenai \emph{topologi lemah}. Untuk pengantar yang lebih terperinci
mengenai kelas topologi ini, lihat buku teks topologi mana pun yang baik atau bagian 11.4 dari catatan saya~\cite{Erdman:2007}.

\begin{defn}\label{C021414} Misalkan $S$ suatu himpunan, untuk setiap $\alpha \in A$ (dengan
$A$ sebarang himpunan indeks) $X_\alpha$ merupakan ruang topologis, dan $f_\alpha
\colon S \sto X_\alpha$ untuk setiap $\alpha \in A$. Misalkan
   \[ \sfml S = \{{f_\alpha}^\gets(U_\alpha) \colon \open{U_\alpha}{X_\alpha}
                                             \text{ dan } \alpha \in A\}\,. \]
Gunakan keluarga $\sfml S$ sebagai subbasis bagi suatu topologi pada~$S$. Topologi ini disebut
 \index{lemah!topologi!diinduksi oleh keluarga fungsi}%
 \index{topologi!lemah!diinduksi oleh keluarga fungsi}%
\df{topologi lemah} yang diinduksi oleh (atau ditentukan oleh) fungsi-fungsi~$f_\alpha$.
\end{defn}

\begin{prop}\label{C021421} Dalam topologi lemah (yang didefinisikan di atas) pada himpunan $S$, setiap fungsi $f_\alpha$ kontinu;
bahkan, topologi lemah adalah topologi paling lemah pada
 \index{lemah!topologi!karakterisasi}%
 \index{topologi!lemah!karakterisasi}%
$S$ yang menjadikan fungsi-fungsi tersebut kontinu.
\end{prop}

\begin{prop}\label{C021424} Misalkan $X$ ruang topologis dengan topologi lemah yang ditentukan
oleh keluarga fungsi $\fml F$. Buktikan bahwa fungsi $g\colon W \sto X$, dengan $W$
ruang topologis, kontinu jika dan hanya jika $f \circ g$ kontinu untuk setiap $f
\in \fml F$.
\end{prop}

\begin{defn}\label{C021437} Misalkan $\bigl(A_\lambda\bigr)_{\lambda \in \Lambda}$ keluarga terindeks
ruang topologis. Topologi lemah pada hasil kali Kartesius $\prod
A_\lambda$ yang diinduksi oleh keluarga proyeksi koordinat $\pi_\lambda$ (lihat
definisi~\ref{C015531}) disebut
 \index{hasil kali!topologi}%
 \index{topologi!hasil kali}%
\df{topologi hasil kali}.
\end{defn}

\begin{prop} Hasil kali dari suatu keluarga ruang Hausdorff adalah Hausdorff.
\end{prop}

\begin{exam}\label{C031473} Misalkan $Y$ ruang topologis dengan topologi lemah (lihat
definisi~\ref{C021414}) yang ditentukan oleh keluarga fungsi terindeks $(f_\alpha)_{\alpha \in A}$,
dengan sifat bahwa, untuk setiap $\alpha \in A$, kodomain $f_\alpha$ berupa ruang topologis~$X_\alpha$.
Maka jaring $\bigl(y_{{}_\lambda}\bigr)_{\lambda \in \Lambda}$ dalam $Y$
konvergen ke titik $a \in Y$ jika dan hanya jika
\mbox{$f_\alpha(y_{{}_\lambda})~\to_{\lambda \in \Lambda}~f_\alpha(a)$} untuk setiap $\alpha
\in A$.
\end{exam}

\begin{exam}\label{C031475} Jika $Y = \prod_{\alpha \in A} X_\alpha$ adalah hasil kali dari ruang-ruang topologis tak kosong
dengan topologi hasil kali (lihat definisi~\ref{C021437}), maka suatu jaring
$\bigl(\sbsb y\lambda\bigr)$ dalam $Y$ konvergen ke $a \in Y$ jika dan hanya jika $\bigl(\sbsb
y\lambda\bigr)_\alpha \sto a_\alpha$ untuk setiap $\alpha \in A$.
\end{exam}

\begin{exam}\label{exam_prod_top_norm} Jika $V$ dan $W$ ruang linear bernorma, maka topologi hasil kali pada
$V \times W$ sama dengan topologi yang diinduksi oleh norma hasil kali pada $V \oplus W$.
\end{exam}

Contoh berikut menggambarkan ketidakcukupan barisan ketika kita menangani ruang topologis umum.

\begin{exam} Misalkan $\fml F = \fml F([0,1])$ adalah himpunan semua fungsi $f\colon [0,1] \sto \R$, dan beri
$\fml F$ topologi hasil kali, yaitu topologi lemah yang ditentukan oleh fungsional-fungsional
evaluasi
    \[ E_x\colon \fml F \sto \R\colon f \mapsto f(x) \]
dengan $x \in [0,1]$. Jadi, lingkungan terbuka dasar dari suatu titik $g \in \fml F$ ditentukan
oleh himpunan titik hingga $A \in \fin[0,1]$ dan bilangan $\epsilon > 0$:
  \[ U(g\,;\,A\,;\, \epsilon) := \{f \in \fml F\colon \abs{f(t) - g(t)} < \epsilon
                           \text{ untuk semua $t \in A$}\}. \]
(Apa nama yang biasa digunakan untuk topologi ini?) Misalkan $\fml G$ himpunan semua fungsi dalam $\fml F$
yang bertumpuan hingga. Maka
  \begin{enumerate}
    \item fungsi konstan $\vc 1$ termasuk dalam $\clo{\fml G}$,
    \item terdapat jaring fungsi dalam $\fml G$ yang konvergen ke~$\vc 1$, tetapi
    \item tidak ada barisan dalam $\fml G$ yang konvergen ke~$\vc 1$.
  \end{enumerate}
\end{exam}

\begin{defn} Jaring $(x_{{}_{\sst\lambda}})$ dalam ruang Hilbert $H$ dikatakan
 \index{konvergensi!dalam ruang Hilbert!lemah}%
 \index{lemah!konvergensi!dalam ruang Hilbert}%
\df{konvergen secara lemah} ke titik $a$ dalam $H$ jika $\langle x{{}_{\sst\lambda}}, y \rangle \to \langle a,y \rangle$
untuk setiap $y \in H$. Dalam hal ini kita menulis $x_{{}_{\sst\lambda}} \to^w a$. Dalam ruang Hilbert, suatu jaring dikatakan
 \index{konvergensi!dalam ruang Hilbert!kuat}%
 \index{kuat!konvergensi!dalam ruang Hilbert}%
\df{konvergen secara kuat} (atau
 \index{konvergensi!dalam norma}%
 \index{norma!konvergensi}%
\df{konvergen dalam norma}) jika jaring tersebut konvergen terhadap topologi yang diinduksi oleh norma. Inilah jenis
konvergensi yang biasa dalam ruang Hilbert. Jika kita ingin menekankan perbedaan di antara cara-cara konvergensi, kita dapat menulis
$x_{{}_{\sst\lambda}} \to^s a$ untuk konvergensi kuat.
\end{defn}

\begin{exer} Jelaskan mengapa topologi pada ruang Hilbert yang dibangkitkan oleh konvergensi lemah jaring sesungguhnya merupakan
topologi lemah dalam pengertian topologis yang lazim.
\end{exer}

Ketika kita mengatakan bahwa suatu subhimpunan dari ruang Hilbert $H$ bersifat
 \index{lemah!ketertutupan}%
 \index{tertutup!secara lemah}%
\df{tertutup secara lemah}, maksudnya subhimpunan tersebut tertutup terhadap topologi lemah pada $H$; suatu himpunan bersifat
 \index{lemah!kekompakan}%
 \index{kompak!secara lemah}%
\df{kompak secara lemah} jika kompak dalam topologi lemah pada~$H$; dan seterusnya. Suatu fungsi
$f\colon H \sto H$ bersifat
 \index{lemah!kontinuitas}%
 \index{kontinu!secara lemah}%
\df{kontinu secara lemah} jika kontinu sebagai pemetaan dari $H$ yang dilengkapi topologi lemah ke $H$
yang dilengkapi topologi lemah.

\begin{exer} Misalkan $H$ ruang Hilbert dan $(x_{{}_{\sst\lambda}})$ jaring dalam $H$.
 \begin{enumerate}
  \item Jika $x_{{}_{\sst\lambda}} \to^s  a$ dalam $H$, maka $x_{{}_{\sst\lambda}} \to^w a$ dalam $H$.
  \item Topologi kuat (norma) pada $H$ lebih kuat (lebih besar) daripada topologi lemah.
  \item Tunjukkan dengan contoh bahwa kebalikan dari (a) tidak berlaku.
  \item Apakah norma pada $H$ kontinu secara lemah?
  \item Jika $x_{{}_{\sst\lambda}} \to^w  a$ dalam $H$ dan jika $\norm{x_{{}_{\sst\lambda}}} \to \norm{a}$, maka
$x_{{}_{\sst\lambda}}\to^s a$ dalam~$H$.
  \item Jika pemetaan linear $T\colon H \sto H$ kontinu (secara kuat), maka pemetaan tersebut kontinu secara lemah.
 \end{enumerate}
\end{exer}

\begin{exer} Misalkan $H$ ruang Hilbert berdimensi tak hingga. Tinjau subhimpunan-subhimpunan berikut dari~$H$:
 \begin{enumerate}[label=\rm{(\alph*)}]
  \item bola satuan terbuka;
  \item bola satuan tertutup;
  \item sfer satuan;
  \item subruang linear tertutup.
 \end{enumerate}
Tentukan untuk setiap himpunan tersebut apakah masing-masing bersifat: tertutup secara kuat, tertutup secara lemah,
terbuka secara kuat, terbuka secara lemah, kompak secara kuat, kompak secara lemah. (Salah satu bagian
terlalu sulit untuk saat ini: \emph{Teorema Alaoglu}~\futurexref{6.2.9}{C067441} akan menunjukkan bahwa
bola satuan tertutup kompak secara lemah.)
\end{exer}

\begin{exer} Misalkan $\{e^n\colon n \in \N\}$ basis biasa untuk $l_2$. Untuk
setiap $n \in \N$, misalkan $a_n = \sqrt n\, e^n$. Manakah di antara pernyataan berikut
yang benar?
 \begin{enumerate}
  \item $0$ adalah titik akumulasi lemah dari himpunan
$\{a_n\colon n \in \N\}$.
  \item $0$ adalah titik gugus lemah dari barisan~$(a_n)$.
  \item $0$ adalah limit lemah dari suatu subbarisan dari~$(a_n)$.
 \end{enumerate}
\end{exer}










\section{Morfisme Universal}
Banyak bagian matematika melibatkan pembentukan objek baru dari objek yang sudah ada---misalnya
hasil kali, koproduk, hasil bagi, pelengkapan, kompaktifikasi, dan unitalisasi.
Sering kali suatu konstruksi dapat---dan sangat baik jika dapat---dicirikan melalui
sebuah diagram yang menjelaskan apa yang “dilakukan” oleh objek hasil konstruksi, alih-alih
menjelaskan apa objek itu atau bagaimana objek tersebut dibangun. Diagram semacam ini disebut
 \index{universal!pemetaan!diagram}%
 \index{pemetaan!universal}%
\df{diagram pemetaan universal}, dan diagram itu menjelaskan
 \index{universal!pemetaan!sifat}%
 \index{pemetaan!diagram!universal}%
\df{sifat universal} objek yang dibangun.
Secara khusus, suatu konstruksi biasanya dapat dicirikan oleh adanya
sebuah morfisme tunggal yang memiliki sifat tertentu. Karena morfisme ini beserta
sifat yang bersesuaian dengannya mencirikan secara tunggal konstruksi tersebut, keduanya masing-masing
disebut \emph{morfisme universal} dan \emph{sifat universal}. Definisi berikut
merupakan salah satu cara untuk menjelaskan peran morfisme semacam itu. Jika ini pertama kalinya Anda
menjumpai konsep \emph{universal}, jangan khawatir karena sifatnya memang cukup abstrak. Menurut saya,
tidak banyak orang yang langsung merasa nyaman dengan definisi ini sebelum menjumpai setidaknya
setengah lusin contoh dalam berbagai konteks. (Ada definisi yang bahkan lebih teknis untuk konsep
penting ini, misalnya \emph{objek awal atau objek terminal dalam kategori koma}. Jika tertarik
menelusurinya, carilah “universal property” di Wikipedia~\cite{wiki:xxx}.)

\begin{defn}\label{00078} Misalkan $\cat A \to^{\ftr{F}} \cat B$ suatu funktor antara kategori
$\cat A$ dan $\cat B$, dan misalkan $B$ suatu objek di~$\cat B$. Pasangan $(A,u)$, dengan $A$
suatu objek di $\cat A$ dan $u$ suatu morfisme di $\cat B$ dari $B$ ke~$\ftr F(A)$, disebut
 \index{universal!morfisme}%
 \index{morfisme!universal}%
\df{morfisme universal} untuk $B$ (terhadap funktor~$\ftr F$) jika untuk setiap
objek $A'$ di $\cat A$ dan setiap morfisme $B \to^f \ftr F(A')$ di $\cat B$ terdapat
sebuah morfisme tunggal $A \to^{\tilde f} A'$ di $\cat A$ sehingga diagram berikut komutatif.
 \begin{equation}\label{00078i}
   \xy
     \qtriangle/{>}`{>}`{-->}/[B`\ftr F(A)`\ftr F(A');u`f`\ftr F(\tilde f)]
     \morphism(1000,500)|r|/{-->}/<0,-500>[A`A';\tilde f]
   \endxy
 \end{equation}
Dalam konteks ini, objek $A$ sering disebut
 \index{universal!objek}%
 \index{objek!universal}%
\df{objek universal} di~$\cat A$.
\end{defn}

Berikut adalah contoh pertama sifat semacam itu.

\begin{defn} Misalkan $F$ suatu objek dalam kategori konkret $\cat C$, dan
$\iota\colon S \sto \abs F$ suatu pemetaan inklusi yang domainnya merupakan himpunan tak kosong~$S$.
Kita mengatakan bahwa objek $F$
 \index{bebas!objek}%
 \index{objek!bebas}%
\df{bebas pada} himpunan $S$ (atau bahwa $F$ adalah \df{objek bebas yang dibangkitkan oleh}~$S$)
jika untuk setiap objek $A$ di $\cat C$ dan setiap pemetaan $f\colon S \sto \abs A$ terdapat
sebuah morfisme tunggal $\widetilde f_\iota\colon F \sto A$ di $\cat C$ sehingga
$\abs{\wt f_\iota} \circ \iota = f$.
 \[\xy
     \qtriangle/{>}`{>}`{-->}/[S`\abs F`\abs A;\iota`f`\abs{\wt f_\iota}]
     \morphism(1000,500)|r|/{-->}/<0,-500>[F`A;\widetilde f_\iota]
   \endxy\]
Kita akan menaruh perhatian pada
 \index{bebas!ruang vektor}%
 \index{ruang!vektor!bebas}%
\df{ruang vektor bebas}, yaitu objek bebas dalam kategori $\cat{VEC}$ yang objeknya adalah
ruang vektor dan morfismenya adalah pemetaan linear. Tentu saja, sekadar
\emph{mendefinisikan} suatu konsep tidak menjamin keberadaannya. Ternyata ruang vektor bebas
memang ada pada setiap himpunan. (Lihat contoh~\ref{ump001z}.)
\end{defn}

\vskip 5 pt

\begin{exer} Dalam definisi sebelumnya, rujukan kepada funktor pelupa sering
dihilangkan dan diagram yang menyertainya sering digambar sebagai berikut:
 \[\xy
     \qtriangle/{>}`{>}`{-->}/[S`F`A;\iota`f`\wt f_\iota]
   \endxy\]
Diagram ini tentu terlihat jauh lebih sederhana. Menurut Anda, mengapa saya memilih versi yang
lebih rumit?
\end{exer}

\vskip 5 pt

\begin{prop} Jika dua objek dalam suatu kategori konkret bebas pada himpunan yang sama,
maka kedua objek tersebut isomorfik.
\end{prop}

\vskip 5 pt

\begin{defn} Misalkan $A$ subhimpunan dari himpunan tak kosong~$S$, dan misalkan $\K$
adalah $\R$ atau~$\C$. Definisikan $\sbsb\chi A \colon S \sto \K$, yaitu
 \index{karakteristik!fungsi}%
 \index{<@$\sbsb{\chi}{A}$ (fungsi karakteristik dari $A$)}%
 \index{fungsi!karakteristik}%
\df{fungsi karakteristik} dari~$A$, dengan
  \[ \sbsb{\chi}{A}(x) = \begin{cases}
                                  1, &\text{jika $x \in A$} \\
                                  0, &\text{selain itu}
                              \end{cases} \]
\end{defn}

\begin{exam}\label{ump001z} Jika $S$ adalah sebarang himpunan tak kosong,
maka terdapat ruang vektor $V$ atas~$\K$ yang bebas pada~$S$. Ruang vektor ini
unik (hingga isomorfisme).
\end{exam}

\begin{proof}[\emph{Petunjuk pembuktian}] Untuk himpunan $S$, ambil $V$ sebagai himpunan
semua fungsi bernilai $\K$ pada $S$ yang bertumpuan hingga. Definisikan penjumlahan dan
perkalian skalar secara titik demi titik. Pemetaan $\iota\colon s \mapsto \sbsb \chi {\{s\}}$
yang membawa setiap unsur $s \in S$ ke fungsi karakteristik dari~$\{s\}$ adalah injeksi
yang diinginkan. Untuk membuktikan bahwa $V$ bebas pada $S$, tunjukkan bahwa untuk setiap
ruang vektor $W$ dan setiap fungsi \hbox{$S~\to^{f}~\abs W$} terdapat pemetaan linear
tunggal $V \to^{\wt f} W$ yang membuat diagram berikut komutatif.
 \[\xy
     \qtriangle/{>}`{>}`{-->}/[S`\abs V`\abs W;\iota`f`\abs{\wt f}]
     \morphism(1100,500)|r|/{-->}/<0,-500>[V`W;\wt f]
   \endxy\]
\ns \end{proof}

\begin{prop} Setiap ruang vektor adalah bebas.
\end{prop}

\begin{proof}[\emph{Petunjuk pembuktian}] Tentu saja, sebagian dari persoalannya adalah
menentukan himpunan $S$ yang membuat ruang vektor tersebut bebas. \ns
\end{proof}


\begin{exam} Misalkan $S$ suatu himpunan tak kosong dan $S' = S \cup \{\vc 1\}$, dengan
$\vc 1$ suatu unsur yang tidak termasuk dalam~$S$. Sebuah
 \index{kata}%
\df{kata} dalam bahasa~$S$ adalah suatu barisan $s$ berunsur di $S'$ yang akhirnya selalu
bernilai $\vc 1$ dan memenuhi syarat: jika $s_k = \vc 1$, maka $s_{k+1} = \vc 1$.
Barisan konstan $(\vc 1, \vc 1, \vc 1, \dots)$ disebut
 \index{kata kosong}%
 \index{kata!kosong}%
\df{kata kosong}. Misalkan $F$ adalah himpunan semua kata dalam bahasa~$S$. Andaikan
$s = (s_1, \dots, s_m, \vc 1, \vc 1, \dots)$ dan
$t = (t_1, \dots, t_n, \vc 1, \vc 1, \dots)$ adalah kata (dengan
$s_1$, \dots, $s_m$, $t_1$, \dots, $t_n \in S$). Definisikan
   \[ s \ast t := (s_1, \dots, s_m, t_1, \dots, t_n, \vc 1, \vc 1, \dots). \]
Operasi ini disebut
 \index{konkatenasi}%
\df{konkatenasi}. Mudah dilihat bahwa himpunan $F$ dengan operasi konkatenasi merupakan
 \index{monoid}%
monoid (semigrup dengan unsur identitas), dengan kata kosong sebagai unsur identitas. Inilah
 \index{bebas!monoid}%
 \index{monoid!bebas}%
\df{monoid bebas} yang dibangkitkan oleh~$S$. Jika kata kosong dikeluarkan, kita memperoleh
 \index{bebas!semigrup}%
 \index{semigrup!bebas}%
\df{semigrup bebas} yang dibangkitkan oleh~$S$. Diagram terkait adalah
 \[\xy
     \qtriangle/{>}`{>}`{-->}/[S`\abs F`\abs G;\iota`f`\abs{\tilde f}]
     \morphism(1100,500)|r|/{-->}/<0,-500>[F`G;\tilde f]
   \endxy\]
di mana $\iota$ adalah injeksi yang jelas $s \mapsto (s,\vc1, \vc 1, \dots)$
(biasanya diperlakukan sebagai pemetaan inklusi), $G$ adalah sebarang semigrup,
$f\colon S \sto G$ adalah sebarang fungsi, dan $\abs{\hphantom{O}}$ adalah funktor pelupa
dari kategori monoid dan homomorfisme (atau kategori semigrup dan homomorfisme) ke kategori
$\cat{SET}$.
\end{exam}

\begin{exam} Ingat kembali penyajian biasa untuk koproduk dua objek dalam suatu kategori
dari definisi~\ref{C015611}. Jika $A_1$ dan $A_2$ adalah dua objek dalam kategori $\cat C$,
maka sebuah
 \index{koproduk}%
\df{koproduk} dari $A_1$ dan $A_2$ adalah tripel $(Q,\iota_1,\iota_2)$, dengan $Q$ suatu
objek di $\cat C$ dan $A_k \to^{\iota_k} Q$ ($k = 1,2$) morfisme di $\cat C$, yang memenuhi
syarat berikut: jika $B$ adalah sebarang objek di $\cat C$ dan
$A_k \to^{f_k} B$ ($k = 1$,$2$) adalah sebarang morfisme di~$\cat C$, maka terdapat sebuah
morfisme tunggal $Q \to^f B$ di $\cat C$ yang membuat diagram berikut komutatif.
  \begin{equation}
    \xy
      \Atrianglepair/<-`<-`<-`>`<-/[B`A_1`Q`A_2;f_1`f`f_2`\iota_1`\iota_2]
    \endxy
   \end{equation}

Sekilas mungkin tidak jelas bahwa konstruksi ini bersifat \emph{universal} dalam arti
definisi~\ref{00078}. Untuk melihatnya, misalkan $\ftr D$ adalah funktor diagonal dari
kategori $\cat C$ ke kategori pasangan $\cat{C^2}$ (lihat contoh~\ref{0007606}).
Andaikan $(Q,\iota_1,\iota_2)$ adalah koproduk objek $A_1$ dan $A_2$ di $\cat C$ dalam
pengertian di atas. Maka $A = (A_1,A_2)$ adalah objek di $\cat{C^2}$,
$A \to^{\iota} \ftr D(Q)$ adalah morfisme di $\cat{C^2}$, dan pasangan
$(Q,\iota)$ bersifat universal dalam arti~\ref{00078}. Diagram yang bersesuaian dengan
diagram~\eqref{00078i} adalah
 \begin{equation}
   \xy
     \qtriangle/{>}`{>}`{-->}/[A`\ftr D(Q)`\ftr D(B);\iota`f`\ftr D(\tilde f)]
     \morphism(1000,500)|r|/{-->}/<0,-500>[Q`B;\tilde f]
   \endxy
 \end{equation}
di mana $B$ adalah sebarang objek di $\cat C$ dan, untuk $k = 1,2$,
$A_k \to^{f_k} B$ adalah sebarang morfisme di $\cat C$, sehingga
$f = (f_1,f_2)$ merupakan morfisme di $\cat{C^2}$.
\end{exam}

\begin{exam} Koproduk dua objek $H$ dan $K$ dalam kategori $\cat{HIL}$, yaitu kategori
 \index{koproduk!dalam $\cat{HIL}$}%
 \index{hil@$\cat{HIL}$!koproduk dalam}%
ruang Hilbert dan pemetaan linear terbatas (dan, secara lebih umum, dalam kategori ruang
hasil kali dalam dan pemetaan linear), adalah jumlah langsung ortogonal eksternalnya
$H \oplus K$ (lihat~\ref{0001504}).
\end{exam}


\begin{exam}\label{00078036} Misalkan $A$ dan $B$ ruang Banach. Pada hasil kali Kartesius
 \index{koproduk!dalam $\cat{BAN_\infty}$}%
 \index{ban@$\cat{BAN_\infty}$!koproduk dalam}%
$A \times B$, definisikan penjumlahan dan perkalian skalar secara titik demi titik. Untuk
setiap $(a,b) \in A \times B$, tetapkan $\norm{(a,b)} = \max\{\norm a,\norm b\}$. Dengan
demikian, $A \times B$ menjadi ruang Banach yang dinotasikan dengan $A \oplus B$ dan
disebut
 \index{jumlah!langsung!ruang Banach}%
\df{jumlah langsung} dari $A$ dan~$B$. Jumlah langsung ini merupakan koproduk dalam
kategori topologis $\cat{BAN_\infty}$ dari ruang Banach, tetapi bukan dalam kategori
geometris yang bersesuaian, yaitu~$\cat{BAN_1}$.
\end{exam}

\begin{exam}\label{00078037} Untuk membangun koproduk dalam kategori geometris
 \index{koproduk!dalam $\cat{BAN_1}$}%
 \index{ban@$\cat{BAN_1}$!koproduk dalam}%
$\cat{BAN_1}$ dari ruang Banach, definisikan operasi ruang vektor secara titik demi titik
pada $A \times B$, tetapi gunakan norma
$\norm{(a,b)}_1 = \norm a + \norm b$ untuk setiap $(a,b) \in A \times B$.
\end{exam}

Hampir setiap konsep dalam teori kategori memiliki konsep dual---konsep yang diperoleh
dengan membalik arah semua panah. Sebagai contoh, kita dapat membalik semua panah dalam
diagram~\eqref{00078i}.


\begin{defn}\label{000781} Misalkan $\cat A \to^{\ftr{F}} \cat B$ suatu funktor antara
kategori $\cat A$ dan $\cat B$, dan misalkan $B$ suatu objek di~$\cat B$. Pasangan $(A,u)$,
dengan $A$ suatu objek di $\cat A$ dan $u$ suatu morfisme di $\cat B$ dari $\ftr F(A)$
ke~$B$, disebut
 \index{kouniversal!morfisme}%
 \index{morfisme!kouniversal}%
\df{morfisme kouniversal untuk $B$ (terhadap~$\ftr F$)} jika untuk setiap objek $A'$ di
$\cat A$ dan setiap morfisme $\ftr F(A') \to^f B$ di $\cat B$ terdapat sebuah morfisme
tunggal $A' \to^{\tilde f} A$ di $\cat A$ sehingga diagram berikut komutatif.
 \begin{equation}\label{000781i}
   \xy
     \qtriangle/{<-}`{<-}`{<--}/[B`\ftr F(A)`\ftr F(A');u`f`\ftr F(\tilde f)]
     \morphism(1000,500)|r|/{<--}/<0,-500>[A`A';\tilde f]
   \endxy
 \end{equation}
Sebagian penulis membalik konvensi ini dan menyebut morfisme dalam~\ref{00078}
\emph{kouniversal}, sedangkan morfisme yang didefinisikan di sini disebut
\emph{universal}. Penulis lain, termasuk penulis buku ini, menyebut keduanya
\emph{morfisme universal}.
\end{defn}

\begin{conv} Morfisme dari kedua jenis yang didefinisikan dalam~\ref{00078}
dan~\ref{000781} akan disebut
 \index{konvensi!tentang universalitas}%
\emph{morfisme universal}.
\end{conv}

\begin{exam}\label{0007812} Ingat kembali pendekatan kategoris terhadap hasil kali dari
definisi~\ref{C015511}. Jika $A_1$ dan $A_2$ adalah dua objek dalam kategori $\cat C$,
maka sebuah
 \index{hasil kali}%
\df{hasil kali} dari $A_1$ dan $A_2$ adalah tripel $(P,\pi_1,\pi_2)$, dengan $P$ suatu
objek di $\cat C$ dan $P \to^{\pi_k} A_k$ ($k = 1,2$) morfisme di $\cat C$, yang memenuhi
syarat berikut: jika $B$ adalah sebarang objek di $\cat C$ dan
$B \to^{f_k} A_k$ ($k = 1,2$) adalah sebarang morfisme di~$\cat C$, maka terdapat sebuah
morfisme tunggal $B \to^f P$ di $\cat C$ yang membuat diagram berikut komutatif.
  \begin{equation}
    \xy
      \Atrianglepair/>`>`>`<-`>/[B`A_1`P`A_2;f_1`f`f_2`\pi_1`\pi_2]
    \endxy
   \end{equation}
Konstruksi ini bersifat (ko)universal dalam arti definisi~\ref{000781}.
\end{exam}

\begin{exam} Hasil kali dua objek $H$ dan $K$ dalam kategori $\cat{HIL}$, yaitu kategori
 \index{hasil kali!dalam $\cat{HIL}$}%
 \index{hil@$\cat{HIL}$!hasil kali dalam}%
ruang Hilbert dan pemetaan linear terbatas (dan, secara lebih umum, dalam kategori ruang
hasil kali dalam dan pemetaan linear), adalah jumlah langsung ortogonal eksternalnya
$H \oplus K$ (lihat~\ref{0001504}).
\end{exam}


\begin{exam} Jika $A$ dan $B$ adalah aljabar Banach, jumlah langsungnya $A \oplus B$
merupakan
 \index{hasil kali!dalam $\cat{BA_\infty}$}%
 \index{ba@$\cat{BA_\infty}$!hasil kali dalam}%
 \index{hasil kali!dalam $\cat{BA_1}$}%
 \index{ba@$\cat{BA_1}$!hasil kali dalam}%
hasil kali dalam kategori topologis maupun geometris dari aljabar Banach, yaitu
$\cat{BA_\infty}$ dan $\cat{BA_1}$. Bandingkan dengan keadaan yang dibahas dalam
contoh~\ref{00078036} dan~\ref{00078037}.
\end{exam}

Konstruksi universal dalam kategori apa pun menghasilkan objek yang unik hingga isomorfisme.

\begin{prop}\label{C035562} Objek universal dalam suatu kategori
 \index{universal!objek!keunikan}%
pada hakikatnya unik.
\end{prop}













\section{Pelengkapan Ruang Hasil Kali Dalam}
Sedikit tinjauan ulang tentang pelengkapan ruang metrik akan berguna di sini. Kadang-kadang
kita mungkin tergoda untuk menganggap bahwa pelengkapan suatu ruang metrik $M$ adalah ruang
metrik lengkap yang memiliki subhimpunan padat yang homeomorfik dengan $M$. Masalah dengan
pencirian ini adalah bahwa pencirian tersebut dapat menghasilkan dua “pelengkapan” ruang
metrik yang bahkan tidak homeomorfik.

\begin{exam} Misalkan $M$ adalah interval $(0,\infty)$ dengan metrik biasa. Terdapat dua
ruang metrik lengkap $N_1$ dan $N_2$ sehingga
  \begin{enumerate}
    \item $N_1$ dan $N_2$ \emph{tidak} homeomorfik,
    \item $M$ homeomorfik dengan suatu subhimpunan padat dari $N_1$, dan
    \item $M$ homeomorfik dengan suatu subhimpunan padat dari~$N_2$.
  \end{enumerate}
\end{exam}

Berikut adalah definisi yang tepat.

\begin{defn}\label{C035511} Misalkan $M$ dan $N$ ruang metrik. Kita mengatakan bahwa $N$
adalah
 \index{pelengkapan!ruang metrik}%
 \index{metrik!ruang!pelengkapan}%
\df{pelengkapan} dari $M$ jika $N$ lengkap dan $M$ isometrik dengan suatu subhimpunan
padat dari~$N$.
\end{defn}

Dalam konstruksi bilangan real dari bilangan rasional, salah satu pendekatan adalah
memperlakukan himpunan bilangan rasional sebagai ruang metrik dan mendefinisikan himpunan
bilangan real sebagai pelengkapannya. Salah satu cara baku untuk menghasilkan pelengkapan
memakai kelas-kelas ekuivalensi barisan Cauchy bilangan rasional. Teknik ini berlaku pula
untuk ruang metrik: cara elementer untuk melengkapi sebarang ruang metrik dimulai dengan
mendefinisikan relasi ekuivalensi pada himpunan barisan Cauchy yang berunsur di ruang
tersebut. Berikut adalah pendekatan yang sedikit kurang elementer, tetapi jauh lebih sederhana.

\begin{prop} Setiap ruang metrik $M$ isometrik dengan suatu subruang dari~$\fml B(M)$.
\end{prop}

\begin{proof}[\emph{Petunjuk pembuktian}] Jika $(M,d)$ adalah ruang metrik, tetapkan
$a \in M$. Untuk setiap $x \in M$, definisikan $\phi_x\colon M \sto \R$ dengan
$\phi_x(u) = d(x,u) - d(u,a)$. Tunjukkan terlebih dahulu bahwa
$\phi_x \in \fml B(M)$ untuk setiap $x \in M$. Kemudian tunjukkan bahwa
$\phi\colon M \sto \fml B(M)$ adalah isometri. \ns
\end{proof}

\begin{cor}\label{C035521} Setiap ruang metrik mempunyai pelengkapan.
\end{cor}

Proposisi berikut menunjukkan bahwa pelengkapan ruang metrik bersifat universal dalam
arti definisi~\ref{00078}. Kategori yang dimaksud dalam hasil ini adalah kategori semua
ruang metrik dan pemetaan kontinu seragam, beserta subkategorinya yang terdiri atas semua
ruang metrik lengkap dan pemetaan kontinu seragam. Di sini, funktor pelupa
$\abs{\hphantom{O}}$ hanya melupakan kelengkapan objek (dan tidak mengubah morfisme).

\begin{prop}\label{C035557} Misalkan $M$ ruang metrik dan $\clo M$ pelengkapannya. Jika
$N$ ruang metrik lengkap dan $f\colon M \sto N$ kontinu seragam, maka terdapat sebuah
fungsi kontinu seragam tunggal $\clo f\colon \clo M \sto N$ yang membuat diagram berikut
komutatif.
 \[\xy
     \qtriangle/{>}`{>}`{-->}/[M`\abs{\clo M}`\abs N;\iota`f`\abs{\clo f}]
     \morphism(1100,500)|r|/{-->}/<0,-500>[\clo M`N;\clo f]
   \endxy\]
\end{prop}

Akibat berikut dari proposisi~\ref{C035557} dan~\ref{C035562} memungkinkan kita berbicara
tentang \emph{pelengkapan} suatu ruang metrik.

\begin{cor}\label{C035564} Pelengkapan ruang metrik unik (hingga isometri).
\end{cor}

Sekarang, bagaimana dengan pelengkapan ruang hasil kali dalam? Seperti telah kita lihat,
setiap ruang hasil kali dalam adalah ruang metrik, sehingga ruang itu mempunyai pelengkapan
(sebagai ruang metrik). Pertanyaannya: apakah pelengkapan ini merupakan ruang Hilbert?
Jawabannya, dengan senang hati, adalah \emph{ya}.

\begin{thm} Misalkan $V$ ruang hasil kali dalam dan $H$ pelengkapannya sebagai ruang
metrik. Maka terdapat hasil kali dalam pada $H$ yang
  \begin{enumerate}
    \item merupakan perpanjangan dari hasil kali dalam pada~$V$, dan
    \item menginduksi metrik pada~$H$.
  \end{enumerate}
\end{thm}


















\endinput
